\documentclass[10pt,letterpaper]{article}
\usepackage{cornell-notes}

\title{Annex B.02.09.02: 02-09-02-metamodel-related-to-the-specification-of-a-model-kind}
\author{}
\date{}
\setDocTitle{Annex B.02.09.02: 02-09-02-metamodel-related-to-the-specification-of-a-model-kind}
\setDocSubtitle{ISO/IEC/IEEE 42010:2022}
\setDocOwner{ISO 42010 Cornell Notes}
\setDocDate{\today}
\setCornellCollection{ISO/IEC/IEEE 42010:2022 Cornell Notes}
\setCornellUnitType{Annex}
\setCornellUnitNumber{B.02.09.02}
\setCornellUnitTitle{02-09-02-metamodel-related-to-the-specification-of-a-model-kind}
\hypersetup{pdftitle={Annex B.02.09.02: 02-09-02-metamodel-related-to-the-specification-of-a-model-kind},pdfsubject={Cornell notes},pdfauthor={}}

\begin{document}
\maketitle
\begin{CornellOverviewBox}{Overview}
{\Large\bfseries\textcolor{CNavy}{Annex B.2.9.2: Metamodel related to the specification of a model kind}}\\[3pt]
{\small\textbf{Classification:} Informative annex}\\
{\small\textbf{Learning objective:} Use the informative guidance on metamodel related to the specification of a model kind to reinforce the normative and conceptual material without treating the guidance itself as mandatory.}
\end{CornellOverviewBox}
\vspace{3pt}

\begin{CornellNotesTable}
\CornellNoteRow{Purpose / key concept}{A metamodel presents one or more constructs which are the AD elements that comprise the vocabulary of the model kind and its specification.}
\CornellNoteRow{Core details}{\begin{itemize}[leftmargin=*,nosep,topsep=1pt]
\item There are various ways of representing metamodels.
\end{itemize}}
\CornellNoteRow{Interpretive note}{\begin{itemize}[leftmargin=*,nosep,topsep=1pt]
\item When specification of an architecture viewpoint specifies multiple model kinds it can be useful to specify a single architecture viewpoint related metamodel unifying the definition of the model kinds. Furthermore, it is often helpful to use a unified metamodel to express a set of related architecture viewpoints (such as when defining an ADF or ADL).
\end{itemize}}
\CornellNoteRow{Related sections}{5.2.9}
\CornellNoteRow{Active recall}{\begin{itemize}[leftmargin=*,nosep,topsep=1pt]
\item What is the central role of Metamodel related to the specification of a model kind?
\item How does Metamodel related to the specification of a model kind connect to adjacent architecture-description concepts?
\end{itemize}}
\end{CornellNotesTable}


\begin{CornellSummaryBox}{Summary}
A metamodel presents one or more constructs which are the AD elements that comprise the vocabulary of the model kind and its specification. This material is informative guidance and does not itself create a conformance requirement.
\end{CornellSummaryBox}

\end{document}
